
\documentclass{article}
\usepackage{colortbl}
\usepackage{makecell}
\usepackage{multirow}
\usepackage{supertabular}

\begin{document}

\newcounter{utterance}

\centering \large Interaction Transcript for game `referencegame', experiment `random\_grids', episode 10 with reflexion{-}Qwen3.5{-}9B.
\vspace{24pt}

{ \footnotesize  \setcounter{utterance}{1}
\setlength{\tabcolsep}{0pt}
\begin{supertabular}{c@{$\;$}|p{.15\linewidth}@{}p{.15\linewidth}p{.15\linewidth}p{.15\linewidth}p{.15\linewidth}p{.15\linewidth}}
    \# & $\;$A & \multicolumn{4}{c}{Game Master} & $\;\:$B\\
    \hline

    \theutterance \stepcounter{utterance}  
    & & \multicolumn{4}{p{0.6\linewidth}}{
        \cellcolor[rgb]{0.9,0.9,0.9}{
            \makecell[{{p{\linewidth}}}]{
                \texttt{\tiny{[P1$\langle$GM]}}
                \texttt{You have attempted this same task before and did not succeed. These are the plans you wrote after those attempts.} \\
\texttt{Trial \#1: New plan: My previous description was too vague and relied on relative positions (e.g., "top{-}left," "bottom{-}right") that were either ambiguous or shared with the distractor grids, failing to uniquely identify the target. Specifically, I missed the unique pattern of the third row in the target grid (starting with two filled cells followed by three empty cells), which distinguishes it from Distractor 1 (starting with an empty cell) and Distractor 2 (starting with an empty cell). Next time, I will construct a description based on a specific, unique row or column pattern that is absent in both distractors, such as "the grid where the third row begins with two filled cells followed by three empty cells," ensuring the expression is precise and exhaustive enough to rule out the alternatives.} \\
\\ 
\\ 
\texttt{You are given three grids, where each of them is 5 by 5 in size.} \\
\texttt{Grids have empty cells marked with "▢" and filled cells marked with "X".} \\
\texttt{Your task is to generate a referring expression that best describes the target grid while distinguishing it from the two other distractor grids.} \\
\texttt{The first grid is the target grid, and the following two grids are the distractors.} \\
\\ 
\texttt{Target grid:} \\
\\ 
\texttt{X □ X X □} \\
\texttt{□ X □ X X} \\
\texttt{X X □ □ X} \\
\texttt{□ X X □ X} \\
\texttt{X □ X □ □} \\
\\ 
\texttt{Distractor grid 1:} \\
\\ 
\texttt{□ □ X X □} \\
\texttt{X □ □ X X} \\
\texttt{□ X X □ X} \\
\texttt{X X □ □ X} \\
\texttt{□ X X □ □} \\
\\ 
\texttt{Distractor grid 2:} \\
\\ 
\texttt{X X □ X □} \\
\texttt{□ X X □ X} \\
\texttt{X □ □ X □} \\
\texttt{□ X □ □ X} \\
\texttt{X □ X □ X} \\
\\ 
\texttt{Instruction: Describe the target grid.} \\
\texttt{Generate the referring expression starting with the tag "Expression: " for the given target grid. Omit any other text.} \\
            }
        }
    }
    & \\ \\

    \theutterance \stepcounter{utterance}  
    & \multicolumn{4}{p{0.6\linewidth}}{
        \cellcolor[rgb]{0.9,0.9,0.9}{
            \makecell[{{p{\linewidth}}}]{
                \texttt{\tiny{[P1$\rangle$GM]}}
                \texttt{Expression: The grid where the third row begins with two filled cells followed by three empty cells, whereas the other two grids have their third rows starting with an empty cell or a filled cell followed by an empty one.} \\
            }
        }
    }
    & & \\ \\

    \theutterance \stepcounter{utterance}  
    & & & \multicolumn{2}{p{0.3\linewidth}}{
        \cellcolor[rgb]{0.9,0.9,0.9}{
            \makecell[{{p{\linewidth}}}]{
                \texttt{\tiny{[GM$|$GM]}}
                \texttt{Expression: The grid where the third row begins with two filled cells followed by three empty cells, whereas the other two grids have their third rows starting with an empty cell or a filled cell followed by an empty one.} \\
            }
        }
    }
    & & \\ \\

    \theutterance \stepcounter{utterance}  
    & & \multicolumn{4}{p{0.6\linewidth}}{
        \cellcolor[rgb]{0.9,0.9,0.9}{
            \makecell[{{p{\linewidth}}}]{
                \texttt{\tiny{[GM$\rangle$P2]}}
                \texttt{You have attempted this same task before and did not succeed. These are the plans you wrote after those attempts.} \\
\texttt{Trial \#1: New plan: I incorrectly identified the "First" grid because I failed to verify the bottom{-}right corner condition. The "First" grid has an 'X' at the bottom{-}right, which contradicts the requirement for it to be empty. The "Second" grid satisfies all conditions: top{-}left is 'X', top{-}right is '▢', the second row starts with '▢', and the bottom{-}right is '▢'. Next time, I will systematically check each grid against every clause of the expression before making a selection, ensuring the bottom{-}right corner is explicitly empty.} \\
\texttt{Answer: Second} \\
\\ 
\\ 
\texttt{You are given three grids, where each of them is 5 by 5 in size.} \\
\texttt{Grids have empty cells marked with "▢" and filled cells marked with "X".} \\
\texttt{You are also given a referring expression that describes one of the given grids.} \\
\texttt{Your task is to select the grid that matches the given referring expression.} \\
\texttt{Generate only the number (in text) of the grid that matches the given expression by selecting first, second, or third.} \\
\\ 
\texttt{First:} \\
\\ 
\texttt{□ □ X X □} \\
\texttt{X □ □ X X} \\
\texttt{□ X X □ X} \\
\texttt{X X □ □ X} \\
\texttt{□ X X □ □} \\
\\ 
\texttt{Second:} \\
\\ 
\texttt{X □ X X □} \\
\texttt{□ X □ X X} \\
\texttt{X X □ □ X} \\
\texttt{□ X X □ X} \\
\texttt{X □ X □ □} \\
\\ 
\texttt{Third:} \\
\\ 
\texttt{X X □ X □} \\
\texttt{□ X X □ X} \\
\texttt{X □ □ X □} \\
\texttt{□ X □ □ X} \\
\texttt{X □ X □ X} \\
\\ 
\texttt{Expression: The grid where the third row begins with two filled cells followed by three empty cells, whereas the other two grids have their third rows starting with an empty cell or a filled cell followed by an empty one.} \\
\texttt{Question: Which grid does the expression refer to?} \\
\texttt{Start with the tag "Answer: ", followed by your selection. Omit any other text.} \\
            }
        }
    }
    & \\ \\

    \theutterance \stepcounter{utterance}  
    & & & \multicolumn{4}{p{0.6\linewidth}}{
        \cellcolor[rgb]{0.9,0.9,0.9}{
            \makecell[{{p{\linewidth}}}]{
                \texttt{\tiny{[GM$\langle$P2]}}
                \texttt{Answer: First} \\
            }
        }
    }
     \\ \\

    \theutterance \stepcounter{utterance}  
    & & & \multicolumn{2}{p{0.3\linewidth}}{
        \cellcolor[rgb]{0.9,0.9,0.9}{
            \makecell[{{p{\linewidth}}}]{
                \texttt{\tiny{[GM$|$GM]}}
                \texttt{first} \\
            }
        }
    }
    & & \\ \\

    \theutterance \stepcounter{utterance}  
    & \multicolumn{4}{p{0.6\linewidth}}{
        \cellcolor[rgb]{0.95,0.95,0.95}{
            \makecell[{{p{\linewidth}}}]{
                \texttt{\tiny{[P2$\rightarrow$P1]}}
                \texttt{Post{-}episode Reflexion (not sent to GM):} \\
\texttt{New plan: My previous description failed because it relied on a negative distinction ("whereas the other two grids have...") which is logically weak and prone to error if the distractors share similar features. I missed the opportunity to describe a positive, unique feature of the target grid that no other grid possesses. Specifically, I should have identified that the target grid is the **only** one where the **first row contains exactly three filled cells** (positions 1, 3, and 4), whereas Distractor 1 has only two and Distractor 2 has only two. Next time, I will construct the expression by first identifying a unique positive pattern (like the specific count and arrangement of 'X's in a specific row or column) that exists exclusively in the target grid, and then explicitly stating that pattern without relying on comparisons to the distractors.} \\
            }
        }
    }
    & & \\ \\

    \theutterance \stepcounter{utterance}  
    & & & \multicolumn{4}{p{0.6\linewidth}}{
        \cellcolor[rgb]{0.95,0.95,0.95}{
            \makecell[{{p{\linewidth}}}]{
                \texttt{\tiny{[P1$\rightarrow$P2]}}
                \texttt{Post{-}episode Reflexion (not sent to GM):} \\
\texttt{New plan: I incorrectly identified the "First" grid because I failed to verify the specific pattern of the third row. The expression requires the third row to start with two filled cells ("X X") followed by three empty cells. Upon re{-}examining the grids, the "First" grid's third row is "□ X X □ X" (starts with empty), and the "Second" grid's third row is "X X □ □ X" (starts with filled, but followed by two filled, not three empty). The "Third" grid's third row is "X □ □ X □" (starts with filled, but followed by empty). Wait, let me re{-}read the expression carefully: "third row begins with two filled cells followed by three empty cells".} \\
\texttt{{-} First Grid Row 3: □ X X □ X (Starts with empty) {-}> Fail.} \\
\texttt{{-} Second Grid Row 3: X X □ □ X (Starts with X X, but then □ □ X, not three empty) {-}> Fail.} \\
\texttt{{-} Third Grid Row 3: X □ □ X □ (Starts with X, then empty) {-}> Fail.} \\
\\ 
\texttt{Let me re{-}read the grids provided in the prompt text vs my internal visualization.} \\
\texttt{Grid 1 Row 3: □ X X □ X} \\
\texttt{Grid 2 Row 3: X X □ □ X} \\
\texttt{Grid 3 Row 3: X □ □ X □} \\
\\ 
\texttt{The expression says: "third row begins with two filled cells followed by three empty cells".} \\
\texttt{None of the grids seem to match "X X □ □ □".} \\
\texttt{Let me re{-}read the expression's second part: "whereas the other two grids have their third rows starting with an empty cell or a filled cell followed by an empty one."} \\
\texttt{{-} Grid 1 Row 3 starts with empty (□). Matches "starting with an empty cell".} \\
\texttt{{-} Grid 3 Row 3 starts with filled (X) followed by empty (□). Matches "filled cell followed by an empty one".} \\
\texttt{{-} Grid 2 Row 3 starts with X X. This is the outlier.} \\
\texttt{Perhaps the expression implies the pattern is "X X" then something else, or maybe I misread the grid text in the prompt.} \\
\texttt{Let's look at Grid 2 again:} \\
\texttt{Row 1: X □ X X □} \\
\texttt{Row 2: □ X □ X X} \\
\texttt{Row 3: X X □ □ X} \\
\texttt{Row 4: □ X X □ X} \\
\texttt{Row 5: X □ X □ □} \\
\\ 
\texttt{Wait, the expression says "followed by three empty cells". Grid 2 has "X X □ □ X". That's two empty, one filled.} \\
\texttt{Is it possible the prompt text has a typo or I am misinterpreting "empty"? "▢" is empty.} \\
\texttt{Let's re{-}read the expression: "The grid where the third row begins with two filled} \\
            }
        }
    }
     \\ \\

\end{supertabular}
}

\end{document}
